\section{集合的包含}

\begin{definition}{包含}{}
	$x$和$y$出现的公式$\forall z\left(\left(z\in x\right)\implies\left(z\in y\right)\right)$简写为如下一种表达:$x\subseteq y$,$y\supseteq x$,``$x$包含于$y$'',``$y$包含$x$'',``$x$是$y$的子集''.
	
	$\neg\left(x\subseteq y\right)$简写为$x\nsubseteq y$或$y\nsupseteq x$.
\end{definition}

\begin{metatheorem}{代入}{}
	设$T,U,V$是表达式,$x$是字母,则$\left(V|x\right)\left(T\subseteq U\right)$和$\left(V|x\right)T\subseteq\left(V|x\right)U$相同.
\end{metatheorem}

\begin{metatheorem}{从包含构造公式}{}
	若$T$和$U$是项,则$T\subseteq U$是公式.
\end{metatheorem}

\begin{definition}{包含公式}{}
	每一个形如$T\subseteq U$(其中$T$和$U$是公式)的表达式称为一个包含关系.
\end{definition}

\begin{proposition}{包含关系的自反性}{}
	$x\subseteq x$.
\end{proposition}

\begin{proposition}{包含关系的传递性}{}
	$\left(x\subseteq y\wedge y\subseteq z\right)\implies x\subseteq z$.
\end{proposition}
